
\documentclass{article}
\usepackage{colortbl}
\usepackage{makecell}
\usepackage{multirow}
\usepackage{supertabular}

\begin{document}

\newcounter{utterance}

\centering \large Interaction Transcript for game `imagegame', experiment `compact\_grids', episode 8 with NEW\_Reflect\_Agent.
\vspace{24pt}

{ \footnotesize  \setcounter{utterance}{1}
\setlength{\tabcolsep}{0pt}
\begin{supertabular}{c@{$\;$}|p{.15\linewidth}@{}p{.15\linewidth}p{.15\linewidth}p{.15\linewidth}p{.15\linewidth}p{.15\linewidth}}
    \# & $\;$A & \multicolumn{4}{c}{Game Master} & $\;\:$B\\
    \hline

    \theutterance \stepcounter{utterance}  
    & & \multicolumn{4}{p{0.6\linewidth}}{
        \cellcolor[rgb]{0.9,0.9,0.9}{
            \makecell[{{p{\linewidth}}}]{
                \texttt{\tiny{[P1$\langle$GM]}}
                \texttt{Let us do the following task. The goal is to fill an empty grid that looks like this:} \\
\\ 
\texttt{▢ ▢ ▢ ▢ ▢} \\
\texttt{▢ ▢ ▢ ▢ ▢} \\
\texttt{▢ ▢ ▢ ▢ ▢} \\
\texttt{▢ ▢ ▢ ▢ ▢} \\
\texttt{▢ ▢ ▢ ▢ ▢} \\
\\ 
\texttt{A filled grid below is 5 by 5 and can look like this:} \\
\\ 
\texttt{▢ ▢ ▢ ▢ ▢} \\
\texttt{▢ ▢ E ▢ ▢} \\
\texttt{▢ ▢ ▢ ▢ ▢} \\
\texttt{▢ ▢ ▢ ▢ ▢} \\
\texttt{X X X X X} \\
\\ 
\texttt{I want you to describe this grid to me, step by step. You don't need to describe the empty squares, which are denoted with "▢". Only describe the location of letters in the grid. Then you wait for me to say "What is your next command?", and then you continue with the next step. Take the size of the grid into consideration while giving commands. You can give only one command at each turn but it can refer to single or multiple cells, rows, columns or even the whole grid at once. When you have described everything, you say "Command: DONE". Only give the command and omit everything else.} \\
\\ 
\texttt{For the filled grid above, here are the example steps.} \\
\\ 
\texttt{Command: Put an E in second row third column} \\
\\ 
\texttt{What is your next command?} \\
\texttt{Command: Fill the last row with X} \\
\\ 
\texttt{What is your next command?} \\
\texttt{Command: DONE} \\
\\ 
\texttt{Ok. Please do this for the following example, which is a 5 by 5 grid.} \\
\texttt{U ▢ ▢ ▢ U} \\
\texttt{▢ U ▢ U ▢} \\
\texttt{▢ ▢ U ▢ ▢} \\
\texttt{▢ U ▢ U ▢} \\
\texttt{U ▢ ▢ ▢ U} \\
\texttt{What is your next command?} \\
            }
        }
    }
    & \\ \\

    \theutterance \stepcounter{utterance}  
    & \multicolumn{4}{p{0.45\linewidth}}{
        \cellcolor[rgb]{0.95,0.95,0.95}{
            \makecell[{{p{\linewidth}}}]{
                \texttt{\tiny{[P1$|$internal]}}
                \texttt{Reflections from previous episodes:} \\
\texttt{{-}} \\
\texttt{ I succeeded. I described the grid accurately in two commands: first the column of Gs on the left, then the bottom row of Gs. This matched the example's efficient style—grouping cells by rows/columns instead of listing each individually. Next time, I would do the same: prioritize concise, grouped commands to cover multiple cells at once.} \\
\texttt{{-}} \\
\texttt{ I succeeded in describing the grid correctly.  } \\
\texttt{I identified the pattern of H's forming an X shape and described it efficiently in two commands: first the main diagonal, then the anti{-}diagonal, then DONE.  } \\
\texttt{This approach minimized steps by grouping cells logically.  } \\
\texttt{Next time, I would do the same — look for symmetrical patterns to describe multiple cells in one command.} \\
\texttt{{-}} \\
\texttt{ I succeeded. I described the grid step by step, focusing only on the filled cells: first the top row of Z's, then the second row starting from column 2, then the third row from column 3, fourth row from column 4, and finally the bottom{-}right single Z. I gave one command per turn and ended with DONE. This approach was efficient because I grouped cells by rows and avoided mentioning empty squares. Next time, I would do the same: describe filled cells in logical batches (rows/columns) to minimize turns while staying clear.} \\
            }
        }
    }
    & & \\ \\

    \theutterance \stepcounter{utterance}  
    & \multicolumn{4}{p{0.45\linewidth}}{
        \cellcolor[rgb]{0.95,0.95,0.95}{
            \makecell[{{p{\linewidth}}}]{
                \texttt{\tiny{[P1$|$internal]}}
                \texttt{{[}attempt 1 – valid{]}} \\
\texttt{Command: Fill the main diagonal with U} \\
            }
        }
    }
    & & \\ \\

    \theutterance \stepcounter{utterance}  
    & \multicolumn{4}{p{0.6\linewidth}}{
        \cellcolor[rgb]{0.9,0.9,0.9}{
            \makecell[{{p{\linewidth}}}]{
                \texttt{\tiny{[P1$\rangle$GM]}}
                \texttt{Command: Fill the main diagonal with U} \\
            }
        }
    }
    & & \\ \\

    \theutterance \stepcounter{utterance}  
    & & & \multicolumn{2}{p{0.3\linewidth}}{
        \cellcolor[rgb]{0.9,0.9,0.9}{
            \makecell[{{p{\linewidth}}}]{
                \texttt{\tiny{[GM$|$GM]}}
                \texttt{Command: Fill the main diagonal with U} \\
            }
        }
    }
    & & \\ \\

    \theutterance \stepcounter{utterance}  
    & & \multicolumn{4}{p{0.6\linewidth}}{
        \cellcolor[rgb]{0.9,0.9,0.9}{
            \makecell[{{p{\linewidth}}}]{
                \texttt{\tiny{[GM$\rangle$P2]}}
                \texttt{Let us draw something together. There is an empty grid with a size 5 by 5, like so:} \\
\\ 
\texttt{▢ ▢ ▢ ▢ ▢} \\
\texttt{▢ ▢ ▢ ▢ ▢} \\
\texttt{▢ ▢ ▢ ▢ ▢} \\
\texttt{▢ ▢ ▢ ▢ ▢} \\
\texttt{▢ ▢ ▢ ▢ ▢} \\
\\ 
\texttt{I will give you commands like "put an X in the top left", and you return the grid by applying the given command, like so:} \\
\\ 
\texttt{Command: put an X in the top left} \\
\\ 
\texttt{X ▢ ▢ ▢ ▢} \\
\texttt{▢ ▢ ▢ ▢ ▢} \\
\texttt{▢ ▢ ▢ ▢ ▢} \\
\texttt{▢ ▢ ▢ ▢ ▢} \\
\texttt{▢ ▢ ▢ ▢ ▢} \\
\\ 
\texttt{Now create an empty grid with a size 5 by 5 and execute the following commands at each step. Once you execute the command, return only the grid and exclude all other text from the output.} \\
\\ 
\texttt{Command: Fill the main diagonal with U} \\
            }
        }
    }
    & \\ \\

    \theutterance \stepcounter{utterance}  
    & & & \multicolumn{4}{p{0.45\linewidth}}{
        \cellcolor[rgb]{0.95,0.95,0.95}{
            \makecell[{{p{\linewidth}}}]{
                \texttt{\tiny{[P2$|$internal]}}
                \texttt{Reflections from previous episodes:} \\
\texttt{{-}} \\
\texttt{ I succeeded. I correctly interpreted "first column" as the leftmost vertical column and filled all five cells in that column with G. The grid output was accurate. Next time, I would do the same: interpret column/row descriptions literally and apply the fill consistently across the entire specified range.} \\
\texttt{{-}} \\
\texttt{ Failed. I guessed "a person holding a phone" but the target was "a person taking a selfie." I correctly identified the person and phone but missed the specific action. Next time, I should focus more on the activity or verb in the image, not just the objects.} \\
\texttt{{-}} \\
\texttt{ **Result:** Failed. I didn't guess the secret word "cactus" in time.} \\
\\ 
\texttt{**Key Actions:**  } \\
\texttt{{-} I guessed "plant" (too broad).  } \\
\texttt{{-} I guessed "aloe" (wrong specific).  } \\
\texttt{{-} I ran out of turns before narrowing it down.} \\
\\ 
\texttt{**Next Time:**  } \\
\texttt{Start with broader category guesses ("plant"), then quickly ask about specific features ("desert?", "has spines?", "green and thick?"). Use yes/no questions earlier to narrow faster.} \\
            }
        }
    }
     \\ \\

    \theutterance \stepcounter{utterance}  
    & & & \multicolumn{4}{p{0.45\linewidth}}{
        \cellcolor[rgb]{0.95,0.95,0.95}{
            \makecell[{{p{\linewidth}}}]{
                \texttt{\tiny{[P2$|$internal]}}
                \texttt{{[}attempt 1 – valid{]}} \\
\texttt{U ▢ ▢ ▢ ▢  } \\
\texttt{▢ U ▢ ▢ ▢  } \\
\texttt{▢ ▢ U ▢ ▢  } \\
\texttt{▢ ▢ ▢ U ▢  } \\
\texttt{▢ ▢ ▢ ▢ U} \\
            }
        }
    }
     \\ \\

    \theutterance \stepcounter{utterance}  
    & & & \multicolumn{4}{p{0.6\linewidth}}{
        \cellcolor[rgb]{0.9,0.9,0.9}{
            \makecell[{{p{\linewidth}}}]{
                \texttt{\tiny{[GM$\langle$P2]}}
                \texttt{U ▢ ▢ ▢ ▢  } \\
\texttt{▢ U ▢ ▢ ▢  } \\
\texttt{▢ ▢ U ▢ ▢  } \\
\texttt{▢ ▢ ▢ U ▢  } \\
\texttt{▢ ▢ ▢ ▢ U} \\
            }
        }
    }
     \\ \\

    \theutterance \stepcounter{utterance}  
    & & & \multicolumn{2}{p{0.3\linewidth}}{
        \cellcolor[rgb]{0.9,0.9,0.9}{
            \makecell[{{p{\linewidth}}}]{
                \texttt{\tiny{[GM$|$GM]}}
                \texttt{Invalid grid format} \\
            }
        }
    }
    & & \\ \\

\end{supertabular}
}

\end{document}
